\section*{References}
\small
\raggedright
\setlength{\parindent}{0pt}
\setlength{\parskip}{4pt}
\everypar{\hangindent=1.5em\hangafter=1}

Allcott, H., and Taubinsky, D. (2015). Evaluating behaviorally motivated policy: Experimental evidence from the lightbulb market. \emph{American Economic Review} 105(8), 2501--2538.

Barberis, N., and Xiong, W. (2009). What drives the disposition effect? An analysis of a long-standing preference-based explanation. \emph{Journal of Finance} 64(2), 751--784.

Bernheim, B. D., and Rangel, A. (2009). Beyond revealed preference: Choice-theoretic foundations for behavioral welfare economics. \emph{Quarterly Journal of Economics} 124(1), 51--104.

Blanchflower, D. G., and Oswald, A. J. (2020). Trends in extreme distress in the United States, 1993--2019. \emph{American Journal of Public Health} 110(10), 1538--1544.

Bloomberg Law (2024). Unhappy workers cost US {\char36}1.9 trillion in lost productivity. January 23, 2024.

Bureau of Economic Analysis (2026). Real gross domestic product per capita [A939RX0Q048SBEA]. Retrieved from FRED, Federal Reserve Bank of St.\ Louis, October 2026.

Bureau of Labor Statistics (2026a). Employed full time: Median usual weekly real earnings, wage and salary workers, 16 years and over [LEU0252881600A]. Retrieved from FRED, Federal Reserve Bank of St.\ Louis, October 2026.

Bureau of Labor Statistics (2026b). Employment-population ratio, 25--54 years [LNS12300060], September 2026. Retrieved from FRED, Federal Reserve Bank of St.\ Louis, October 2026.

Cai, Y., and Shefrin, H. (2018). Focal points and firm risk. \emph{European Financial Management} 24(4), 521--544.

Chancelier, J.-Ph., \O ksendal, B., and Sulem, A. (2002). Combined stochastic control and optimal stopping, and application to numerical approximation of combined stochastic and impulse control. \emph{Proceedings of the Steklov Institute of Mathematics} 237, 140--163.

Charles, K. K., Hurst, E., and Roussanov, N. (2009). Conspicuous consumption and race. \emph{Quarterly Journal of Economics} 124(2), 425--467.

Chetty, R., Grusky, D., Hell, M., Hendren, N., Manduca, R., and Narang, J. (2017). The fading American dream: Trends in absolute income mobility since 1940. \emph{Science} 356(6336), 398--406.

Crandall, M. G., Ishii, H., and Lions, P.-L. (1992). User's guide to viscosity solutions of second order partial differential equations. \emph{Bulletin of the American Mathematical Society} 27(1), 1--67.

Dalton, P. S., Ghosal, S., and Mani, A. (2016). Poverty and aspirations failure. \emph{Economic Journal} 126(590), 165--188.

Dixit, A. (1989). Entry and exit decisions under uncertainty. \emph{Journal of Political Economy} 97(3), 620--638.

Dixit, A., and Pindyck, R. S. (1994). \emph{Investment under Uncertainty}. Princeton University Press.

Ebert, S., and Strack, P. (2015). Until the bitter end: On prospect theory in a dynamic context. \emph{American Economic Review} 105(4), 1618--1633.

El Karoui, N. (1981). Les aspects probabilistes du contr\^ole stochastique. In \emph{\'Ecole d'\'Et\'e de Probabilit\'es de Saint-Flour IX--1979}, Lecture Notes in Mathematics 876. Springer.

Fleming, W. H., and Soner, H. M. (2006). \emph{Controlled Markov Processes and Viscosity Solutions} (2nd ed.). Springer.

Gallup (2026). \emph{State of the Global Workplace: 2026 Report}. Washington, DC: Gallup.

Genicot, G., and Ray, D. (2017). Aspirations and inequality. \emph{Econometrica} 85(2), 489--519.

Goggins, D. (n.d.). Taking souls [video]. YouTube Shorts. \url{https://youtube.com/shorts/5fkhdbG6-wo}, accessed October 29, 2026.

Harter, J. (2022). Is quiet quitting real? \emph{Gallup}, September 6, 2022.

Harter, J. (2026). U.S. employee engagement declines from 2020 peak. \emph{Gallup}, January 27, 2026.

Helliwell, J. F., Layard, R., Sachs, J. D., De Neve, J.-E., Aknin, L. B., and Wang, S. (Eds.) (2025). \emph{World Happiness Report 2025}. Wellbeing Research Centre, University of Oxford.

Henderson, V. (2012). Prospect theory, liquidation, and the disposition effect. \emph{Management Science} 58(2), 445--460.

Ingersoll, J. E., and Jin, L. J. (2013). Realization utility with reference-dependent preferences. \emph{Review of Financial Studies} 26(3), 723--767.

Joint Center for Housing Studies of Harvard University (2025). \emph{The State of the Nation's Housing 2025}. Cambridge, MA.

Kahneman, D., and Tversky, A. (1979). Prospect theory: An analysis of decision under risk. \emph{Econometrica} 47(2), 263--291.

Karatzas, I., and Wang, H. (2000). Utility maximization with discretionary stopping. \emph{SIAM Journal on Control and Optimization} 39(1), 306--329. doi:10.1137/S0363012998346323.

K\H{o}szegi, B., and Rabin, M. (2006). A model of reference-dependent preferences. \emph{Quarterly Journal of Economics} 121(4), 1133--1165.

Krylov, N. V. (1980). \emph{Controlled Diffusion Processes}. Springer.

Kyle, A. S., Ou-Yang, H., and Xiong, W. (2006). Prospect theory and liquidation decisions. \emph{Journal of Economic Theory} 129(1), 273--288.

Lee, D., Park, J., and Shin, Y. (2023). Where are the workers? From great resignation to quiet quitting. NBER Working Paper 30833.

Lester, R. (2026). Replication data for: Rational Foreclosure. Harvard Dataverse. \url{https://doi.org/10.7910/DVN/V5NSZX}.

Luttmer, E. F. P. (2005). Neighbors as negatives: Relative earnings and well-being. \emph{Quarterly Journal of Economics} 120(3), 963--1002.

Lybbert, T. J., and Wydick, B. (2018). Poverty, aspirations, and the economics of hope. \emph{Economic Development and Cultural Change} 66(4), 709--753.

March, J. G., and Shapira, Z. (1992). Variable risk preferences and the focus of attention. \emph{Psychological Review} 99, 172--183.

McDonald, R., and Siegel, D. (1986). The value of waiting to invest. \emph{Quarterly Journal of Economics} 101(4), 707--727.

McKean, H. P. (1965). A free boundary problem for the heat equation arising from a problem in mathematical economics. Appendix to a paper by P. A. Samuelson. \emph{Industrial Management Review} 6(2), 32--39.

Miller, G. E., and Wrosch, C. (2007). You've gotta know when to fold 'em: Goal disengagement and systemic inflammation in adolescence. \emph{Psychological Science} 18(9), 773--777.

Murthy, V. H. (2016). Facing addiction in the United States: The Surgeon General's report on alcohol, drugs, and health. \emph{JAMA}, published online November 17, 2016. doi:10.1001/jama.2016.18215.

National Association of Realtors (2025). First-time home buyer share falls to historic low of 21\%, median age rises to 40. News release on the \emph{2025 Profile of Home Buyers and Sellers}, November 4, 2025.

Nesse, R. M. (2000). Is depression an adaptation? \emph{Archives of General Psychiatry} 57(1), 14--20.

Pega, F., N\'afr\'adi, B., Momen, N. C., Ujita, Y., Streicher, K. N., Pr\"uss-\"Ust\"un, A. M., et al. (2021). Global, regional, and national burdens of ischemic heart disease and stroke attributable to exposure to long working hours for 194 countries, 2000--2016: A systematic analysis from the WHO/ILO Joint Estimates of the Work-related Burden of Disease and Injury. \emph{Environment International} 154, 106595.

Peskir, G. (2005). A change-of-variable formula with local time on curves. \emph{Journal of Theoretical Probability} 18(3), 499--535.

Pew Research Center (2025). Views of children's financial future. January 9, 2025.

Pham, H. (2009). \emph{Continuous-time Stochastic Control and Optimization with Financial Applications}. Springer.

Ray, D. (1998). \emph{Development Economics}. Princeton University Press.

Ray, D. (2006). Aspirations, poverty, and economic change. In A. V. Banerjee, R. B\'enabou, and D. Mookherjee (Eds.), \emph{Understanding Poverty}. Oxford University Press.

Richter, C. P. (1957). On the phenomenon of sudden death in animals and man. \emph{Psychosomatic Medicine} 19(3), 191--198.

Shefrin, H. M., and Thaler, R. H. (1988). The behavioral life-cycle hypothesis. \emph{Economic Inquiry} 26(4), 609--643.

Thaler, R. H., and Shefrin, H. M. (1981). An economic theory of self-control. \emph{Journal of Political Economy} 89(2), 392--406.

Thaler, R. H., and Sunstein, C. R. (2008). \emph{Nudge: Improving Decisions about Health, Wealth, and Happiness}. Yale University Press.

Witters, D. (2025). U.S. depression rate remains historically high. \emph{Gallup}, September 9, 2025.

World Health Organization (2019). Burn-out an ``occupational phenomenon'': International Classification of Diseases. News release, May 28, 2019.

World Health Organization (2022). WHO and ILO call for new measures to tackle mental health issues at work. News release, September 28, 2022.

Wrosch, C., Miller, G. E., Scheier, M. F., and Brun de Pontet, S. (2007). Giving up on unattainable goals: Benefits for health? \emph{Personality and Social Psychology Bulletin} 33(2), 251--265.
